\section{Qualitative System Traces}
\label{app:traces}

The following three traces are captured from the live \textbf{EduSHIELD} system deployed on the Programming Fundamentals course. Each trace records: the student query; Stage~1 Boundary Guard verdict with semantic similarity $s_{\text{sem}}$; Stage~2 KG-RAG retrieval with cosine similarities; Stage~3 LLM response text; and Stage~4 Hallucination Guard verdict with positive score $s^+$ (response claim vs.\ KG claims), negation score $s^-$ (response claim vs.\ KG negations), gap $s^{-} - s^{+}$, and the final delivery decision.

\subsection*{D.1\quad Verified In-Domain Response}
The student asks a legitimate in-domain question about for loops. The Boundary Guard passes it (IN\_DOMAIN, $s_{\text{sem}}=0.7132$). The LLM response is grounded in KG-retrieved facts. The Hallucination Guard labels the claim \textsc{Verified} ($s^+ \geq \theta_v = 0.52$, $s^{-} - s^{+} < \delta$). The response is delivered as-is.
\begin{small}
\begin{verbatim}
Student query : "What does a for loop do in python?"

Stage 1 -- Boundary Guard (pre-generation):
  s_sem          = 0.7132
  matched concept= "for loops"
  classification = IN_DOMAIN

Stage 2 -- KG-RAG Retrieval (top-3 unique facts):
  [1] concept = "FOR LOOPS"
      claim   = "A for loop in Python iterates over each item in a sequence"
      sim     = 0.7132

Stage 3 -- LLM Generation:
  "A for loop iterates over a sequence of values a fixed number of times."

Stage 4 -- Hallucination Guard (post-generation):
  claim   = "A for loop iterates over a sequence of values a fixed number of t"
  s+      = 0.5812   (cosine sim: response claim vs KG claims)
  s-      = 0.5419   (cosine sim: response claim vs KG negations)
  s- - s+ = -0.0393  (delta threshold delta = 0.04)
  verdict = VERIFIED

  Overall verdict = VERIFIED

Delivered to student: VERIFIED
  (LLM response delivered as-is)
\end{verbatim}
\end{small}

\subsection*{D.2\quad Contradicted Response (Hallucination Caught)}
A deliberately incorrect claim is injected as the LLM response: the system asserts that a while loop checks its condition \emph{after} the loop body executes (a do-while behaviour, not a while behaviour). The Hallucination Guard finds $s^- > s^+$ by a gap of $+0.1027 > \delta = 0.04$, triggering \textsc{Contradicted}. The original response is replaced by a KG-sourced correction.
\begin{small}
\begin{verbatim}
Student query : "How does a while loop work?"

Stage 1 -- Boundary Guard (pre-generation):
  s_sem          = 0.6833
  matched concept= "while loop"
  classification = IN_DOMAIN

Stage 2 -- KG-RAG Retrieval (top-3 unique facts):
  [1] concept = "WHILE LOOP"
      claim   = "A while loop checks its condition before executing the body"
      sim     = 0.6833

Stage 3 -- LLM Generation:
  "A while loop checks the condition after executing the loop body."

Stage 4 -- Hallucination Guard (post-generation):
  claim   = "A while loop checks the condition after executing the loop body."
  s+      = 0.5744   (cosine sim: response claim vs KG claims)
  s-      = 0.6771   (cosine sim: response claim vs KG negations)
  s- - s+ = +0.1027  (delta threshold delta = 0.04)
  verdict = CONTRADICTED

  Overall verdict = CONTRADICTED

  KG-sourced correction appended:
  "A while loop checks its condition before executing the body.
 If the condition is false initially, the loop body never runs."
  [source: cse1321-loops.pptx]

Delivered to student: CONTRADICTED
  (original LLM claim replaced by KG correction above)
\end{verbatim}
\end{small}

\subsection*{D.3\quad Boundary Redirect (No LLM Call)}
The student asks about recursion, a concept mentioned in course slides 
but reserved for a later course in the sequence. The token-level 
boundary match ($m_{\text{bnd}}(\text{`recursion'}) = 1$) fires before 
the semantic score is evaluated, triggering an immediate \textsc{Boundary} 
verdict ($s_{\text{sem}} = 0.4501$, which would have passed the in-domain 
semantic threshold alone). No LLM call is made; a redirect message is 
delivered directly.
\begin{small}
\begin{verbatim}
Student query : "Can you explain recursion to me?"

Stage 1 -- Boundary Guard (pre-generation):
  s_sem          = 0.4501
  matched concept= "loops"
  classification = BOUNDARY
  m_bnd("recursion") = 1  --> priority override fires
  --> REDIRECT: No LLM call made.

Delivered to student:
  "Recursion is a topic covered in the next course in the sequence.
   Would you like to review for loops or while loops instead?"
\end{verbatim}
\end{small}